\section{Desacoplamiento entre el modelo de planeación y el modelo de ejecución}

\subsection{Por qué se necesita la división Planner-Executor}
En la segunda mitad, el expositor destacó un enfoque jerárquico: el modelo de nivel superior se encarga de planear, y el modelo de nivel inferior se encarga de ejecutar. De esta manera, una tarea compleja puede dividirse en subtareas localmente controlables, lo que reduce el ruido de tomar decisiones directas paso a paso.

\begin{figure}[H]
\centering
\includegraphics[width=0.82\textwidth]{frame-06-planning.jpg}
\caption{Video, aproximadamente 00:27:40: el expositor explica que primero el planning model produce un plan y luego el modelo de ejecución lo concreta en acciones}
\end{figure}

\begin{knowledgebox}{Beneficios del desacoplamiento}
\begin{itemize}
    \item La capa de planeación puede imponer restricciones sobre el objetivo global;
    \item La capa de ejecución puede realizar acciones localmente óptimas para la página actual;
    \item La localización de errores es más clara, lo que facilita la depuración y la reversión.
\end{itemize}
\end{knowledgebox}

\subsection{Cómo diseñar la interfaz jerárquica}
Una interfaz aplicable en la práctica suele contener tres partes: la descripción del subobjetivo, los límites de ejecución y el criterio de finalización. Sin una interfaz clara, la capa de planeación y la capa de ejecución terminan interfiriéndose entre sí, lo que produce el fenómeno de ``el plan está bien escrito pero la ejecución se desvía''.

\begin{importantbox}{Campos mínimos de la interfaz jerárquica}
\begin{enumerate}
    \item \textbf{subgoal}: el objetivo de la etapa actual;
    \item \textbf{allowed actions}: el conjunto de acciones permitidas;
    \item \textbf{stop condition}: el criterio para determinar si se completó o falló;
    \item \textbf{fallback rule}: la condición para regresar a la capa superior en caso de anomalía.
\end{enumerate}
\end{importantbox}

\subsection{Cuándo conviene volver a planear}
En cuanto la ejecución dispare una anomalía de alta confianza (el elemento no existe, un salto de página anómalo, exceso del presupuesto), el sistema debe regresar a tiempo a la capa de planeación, en lugar de continuar ciegamente con la siguiente acción.

\begin{warningbox}{Fallas comunes de la capa de planeación}
\begin{itemize}
    \item La granularidad del plan es demasiado gruesa y la capa de ejecución no puede aplicarlo;
    \item La granularidad del plan es demasiado fina y se pierde el valor de la jerarquía;
    \item Falta la regla de ``cuándo volver a planear'', lo que provoca que el error siga propagándose.
\end{itemize}
\end{warningbox}

\subsection{Resumen de la sección}
El desacoplamiento Planner-Executor es un patrón de ingeniería clave para los agentes de cadena larga. Lo que realmente determina el resultado es el diseño de la interfaz y la estrategia de replaneación, y no solo cambiar a un modelo más grande.

